\documentclass[11pt]{article}
\usepackage{amsmath,amssymb,amsthm,hyperref}
\newtheorem{theorem}{Theorem}
\newtheorem{lemma}[theorem]{Lemma}
\newtheorem{corollary}[theorem]{Corollary}
\title{The exact threshold degree of graph parity on a permutation orbit}
\author{}
\date{November 22, 2026}
\begin{document}
\maketitle

This is an obstruction to a specific uniform permutation decoder, not
a lower bound for the manuscript's Boolean head complexity. In particular,
it does not settle the requested universal upper bound.

Let $p\ge1$, let $a_1<\cdots<a_p$ be real numbers, and let
\[
 \Omega_a=\{(a_{\pi(1)},\ldots,a_{\pi(p)}):\pi\in S_p\}.
\]
Let $E\subseteq\{(i,j):1\le i<j\le p\}$, set $m=|E|$, and put
\[
 Q_E(z)=\prod_{(i,j)\in E}(z_j-z_i),\qquad
 f_E(z)=\operatorname{sign}Q_E(z)\quad(z\in\Omega_a).
\]
The empty product is one. Threshold degree here means the least total
degree of a real polynomial with the specified strict signs on
\emph{every} point of $\Omega_a$.

\begin{lemma}[Divided symmetrization]
For every real polynomial $P$ of total degree less than $m$,
\[
 \sum_{\pi\in S_p}\frac{P(a_{\pi(1)},\ldots,a_{\pi(p)})}
 {Q_E(a_{\pi(1)},\ldots,a_{\pi(p)})}=0.
\]
\end{lemma}
\begin{proof}
Write $\Delta(z)=\prod_{i<j}(z_j-z_i)$ and $D=\binom p2$.
The quotient $R=P\Delta/Q_E$ is a polynomial of degree less than $D$.
For distinct coordinates $z$, the rational sum in the assertion is
\[
 \frac{1}{\Delta(z)}
 \sum_{\pi\in S_p}\operatorname{sign}(\pi)
             R(z_{\pi(1)},\ldots,z_{\pi(p)}).
\]
Its numerator is alternating under every coordinate transposition.
It therefore vanishes on each hyperplane $z_i=z_j$. Each linear
polynomial $z_j-z_i$ divides the numerator. These factors are distinct
nonassociate primes in the real polynomial ring, so their product
$\Delta$ divides it. Its degree is less than $D$, hence the numerator
is zero. Evaluate at $z=a$.
\end{proof}

This cancellation is the graph divided-symmetrization identity discussed
in Petrov, Section~1, pp.~1--2,
\url{https://arxiv.org/pdf/1512.07136}.
Postnikov, Proposition~3.4, gives the path version and the alternating
polynomial argument in
\url{https://math.mit.edu/~apost/papers/permutohedron_full.pdf}.
The proof above includes the identity needed here.

\begin{theorem}[Exact degree on every fixed permutation orbit]
The threshold degree of $f_E$ on $\Omega_a$ is exactly $m$.
More explicitly, an annihilating signed measure for all degrees below
$m$ is
\[
 \mu(z)=\frac1{Q_E(z)}\quad(z\in\Omega_a).
\]
Every weight is nonzero and its sign is $f_E(z)$.
\end{theorem}
\begin{proof}
The polynomial $Q_E$ proves the upper bound. If a polynomial $P$ of
degree less than $m$ has the strict sign $f_E$, then $P(z)/Q_E(z)>0$
at every $z\in\Omega_a$. The sum of these finitely many positive
numbers is positive, contradicting the lemma. When $m=0$, the constant
polynomial one gives degree zero, and no negative degree is needed.
\end{proof}

\begin{corollary}[Cost of a uniform affine-fraction decoder]
Suppose
\[
 S(z)=c+\sum_{h=1}^H\frac{A_h(z)}{B_h(z)},
\]
where $c$ is constant, all $A_h,B_h$ are affine, and every $B_h$ is
positive on $\Omega_a$. If $S$ has sign $f_E$ on every point of
$\Omega_a$, then $H\ge m$.
If an arbitrary free affine term replaces $c$, the conclusion is
$H+1\ge m$.
\end{corollary}
\begin{proof}
Multiply by the positive product of all $B_h$. With a constant term,
the resulting polynomial has degree at most $H$, and its strict
signs are unchanged. Apply the theorem. A free affine term increases
the degree bound to $H+1$.
\end{proof}

\paragraph{Examples.}
A full star has $m=p-1$, so decoding parity of the rank of its center
already requires $p-1$ heads under these hypotheses. The complete graph
has $m=\binom p2$, and its sign is permutation parity. A matching with
$m$ edges gives the usual degree-$m$ parity obstruction, but the theorem
does not require its edges to be disjoint.

\paragraph{Application and exact limitation.}
An affine binary hash of inversion vectors has output bits given, up
to a fixed sign, by $f_E$ for its row edge sets $E$. Thus a proposed
decoder that works uniformly on all permutations of a fixed set of
stored scalar values pays at least $|E|$ affine fractions for that
output bit. Storing many bits by the existence of a surjective hash
does not supply this decoder.

The theorem does not cover a decoder required only on one specially
chosen representative permutation for each hash output. That subset
need not support the annihilating measure. It also does not cover
nonlinear assignment of head coefficients from the stored scalars,
because its affine hypothesis is in $z$. These are concrete limits of
the obstruction, not assumed constructions beyond it.

\paragraph{Audit status.}
This note is a manual proof. It is not a Lean declaration. The script
\texttt{check\_permutation\_parity\_degree.py} independently checks the
signed annihilator on specified finite integer orbits using integer
arithmetic. Those diagnostics are not a substitute for the proof or
evidence for the universal head bound.
\end{document}
